\section{Alignement cyclotomique, orbifold d'ordre \(3\) et Moonshine monstrueux}
\label{p12-83-c20-s8:chap:alineamiento-orbifold-moonshine}

Le réseau de Leech construit dans la
\cref{p12-83-c20-s5:thm:identificacion-vecino-leech} permet de passer de l'incidence
finie à une structure graduée infinie. Cette transition ne s'effectuera pas
au moyen d'une chaîne de noms. Seront construits, dans un ordre typé :
\begin{enumerate}
 \item l'entrelacement des douze fibres \(A_2\) ;
 \item une isométrie de réseau sans points fixes d'ordre trois ;
 \item sa trace graduée et son relèvement de type zéro ;
 \item l'orbifold cyclique et son automorphisme dual de classe \(3B\) ;
 \item l'algèbre d'opérateurs de sommets lunaire ;
 \item par des publications distinctes, son groupe d'automorphismes, son caractère
       gradué et les séries de McKay--Thompson.
\end{enumerate}
Moonshine sera le théorème qui coordonne l'action et les séries. On n'utilisera pas
le raccourci non typé
\(V^\natural\to\mathbb M\to\text{Moonshine}\).

\subsection{Alignement des \(12\) fibres APP et Witt}

Les douze fibres de \(W_{\mathrm{APP}}\cong A_2^{12}\) portent trois
étiquettes indépendantes :
\[
 i\in\mathbb Z/3\mathbb Z
 \quad\text{(paire de coordonnées)},\qquad
 \mu\in\{0,1\}
 \quad\text{(feuillet)},\qquad
 \varepsilon\in\{0,1\}
 \quad\text{(orbite cyclotomique)}.
\]
L'origine dodécaphasique construite dans le TPK étant fixée, nous définissons
\begin{equation}
 \lambda(i,\mu,\varepsilon)
 :=4i+2\mu+\varepsilon\pmod{12}.
 \label{p12-83-c20-s8:eq:alineamiento-doce-fibras-u031}
\end{equation}

\begin{lemma}[Bijectivité de l'alignement]
\label{p12-83-c20-s8:lem:biyectividad-alineamiento-u031}
L'application \(\lambda\) de
\eqref{p12-83-c20-s8:eq:alineamiento-doce-fibras-u031} est une bijection sur
\(\mathbb Z/12\mathbb Z\).
\end{lemma}

\begin{proof}
Pour chaque \(i\), les quatre valeurs \(4i+2\mu+\varepsilon\) forment le bloc
\(\{4i,4i+1,4i+2,4i+3\}\). Les trois blocs pour \(i=0,1,2\) sont disjoints
et couvrent les douze classes.
\end{proof}

Soit \(C\) le Coxeter de \(A_2\) et soit \(\mathsf R\) la réflexion qui
satisfait
\begin{equation}
 \mathsf R C\mathsf R=C^{-1}.
 \label{p12-83-c20-s8:eq:reflexion-conjuga-coxeter-u031}
\end{equation}
Les trivialisations \(\iota_b\) et les cadres orientés \(P_b\) étant fixés,
nous posons
\begin{align}
 g_b&:=\iota_b^{-1}P_bCP_b^{-1}\iota_b,
 \label{p12-83-c20-s8:eq:coxeter-calibrado-fibra-u031}\\
 T_{i\mu\varepsilon}
 &:=\iota_{\lambda(i,\mu,\varepsilon)}^{-1}
 P_{\lambda(i,\mu,\varepsilon)}\mathsf R^\mu.
 \label{p12-83-c20-s8:eq:entrelazador-fibra-u031}
\end{align}

\begin{theorem}[Entrelacement APP--Witt]
\label{p12-83-c20-s8:thm:alineamiento-app-witt-u031}
Pour chaque triplet \((i,\mu,\varepsilon)\),
\begin{equation}
 g_{\lambda(i,\mu,\varepsilon)}T_{i\mu\varepsilon}
 =T_{i\mu\varepsilon}C^{(-1)^\mu}.
 \label{p12-83-c20-s8:eq:identidad-entrelazamiento-app-witt-u031}
\end{equation}
La somme directe définit un isomorphisme \(C_3\)-équivariant
\begin{equation}
 T_\lambda:
 W_{\mathrm{APP}}\otimes\mathbb C
 \xrightarrow{\;\cong\;}
 W_c:=\bigoplus_{b=0}^{11}
 (A_{2,b}\otimes\mathbb C).
 \label{p12-83-c20-s8:eq:isomorfismo-app-witt-u031}
\end{equation}
\end{theorem}

\begin{proof}
Pour \(\mu=0\), la définition de \(g_b\) donne directement la conjugaison
par \(C\). Pour \(\mu=1\), on insère
\(\mathsf RC\mathsf R=C^{-1}\). Le
\cref{p12-83-c20-s8:lem:biyectividad-alineamiento-u031} garantit que la somme directe
utilise exactement une fois chaque position de Witt et conserve paire, feuillet et orbite.
\end{proof}

Le théorème est relatif au transport simultané de l'ordre des paires, des feuillets et des
orbites, de l'origine de Witt, des cadres et de l'orientation. Changer une
seule de ces coordonnées ne définit pas le même foncteur.

\subsection{Trace graduée du générateur ternaire}

Soit
\[
 \rho_W:=\bigoplus_{b=0}^{11}g_b
\]
et considérons l'enveloppe bosonique
\begin{equation}
 \mathcal H_{\rm bos}(W_c)
 :=
 \operatorname{Sym}\!\left(
 \bigoplus_{n\ge1}W_c[n]
 \right).
 \label{p12-83-c20-s8:eq:fock-bosonico-doce-fibras-u031}
\end{equation}
Si \(N\) est l'opérateur de degré, nous définissons
\begin{equation}
 Z_{\rm Fock}(q)
 :=
 q^{-1}\operatorname{Tr}_{\mathcal H_{\rm bos}(W_c)}
 \left(\Gamma_{\rm bos}(\rho_W)q^N\right).
 \label{p12-83-c20-s8:eq:traza-graduada-fock-u031}
\end{equation}

\begin{proposition}[Produit cyclotomique de niveau trois]
\label{p12-83-c20-s8:prop:producto-tres-doce-fibras-u031}
Comme série formelle de Laurent,
\begin{equation}
 \boxed{
 Z_{\rm Fock}(q)
 =q^{-1}\prod_{n\ge1}(1+q^n+q^{2n})^{-12}
 =:t_3(q).}
 \label{p12-83-c20-s8:eq:t3-producto-u031}
\end{equation}
Son développement commence par
\begin{equation}
 t_3(q)
 =q^{-1}-12+54q-76q^2-243q^3+1188q^4+O(q^5).
 \label{p12-83-c20-s8:eq:t3-expansion-u031}
\end{equation}
\end{proposition}

\begin{proof}
Chaque \(g_b\) est conjugué à \(C\), de sorte que
\[
 \det(I-g_bq^n)=1+q^n+q^{2n}.
\]
L’identité
\[
 \sum_{k\ge0}
 \operatorname{Tr}(\operatorname{Sym}^k\rho_W)t^k
 =\det(I-\rho_Wt)^{-1}
\]
appliquée à chaque degré produit le produit. La multiplication formelle jusqu'au
degré cinq donne le développement. En particulier, le coefficient de \(q\)
s'obtient comme
\[
 [q^2]\,
 (1+q+q^2)^{-12}(1+q^2+q^4)^{-12}
 =66-12=54.
\]
\end{proof}

La valeur \(54\) apparaît ici comme coefficient d'une trace construite à partir de
douze fibres. L'unité U030 l'a également obtenue comme indice orienté de
l'agrégat APP et a prouvé que douze est l'unique solution positive de la rigidité
cyclotomique--radiale. L'égalité des deux valeurs est une compatibilité de
deux réalisations ; aucune n'utilise le coefficient de \(J\) comme entrée.

\subsection{Mot de support total et isométrie sans points fixes}

Soit
\[
 q_*=(1,1,1,1,1,1)\in\mathbb F_3^6.
\]
La multiplication par la matrice de Paley--Witt donne
\begin{equation}
 q_*A_W=(2,1,1,1,1,1).
 \label{p12-83-c20-s8:eq:producto-qstar-aw-u031}
\end{equation}
Par conséquent,
\begin{equation}
 c_*:=(q_*,q_*A_W)
 =(1,1,1,1,1,1,2,1,1,1,1,1)
 \in C_W
 \label{p12-83-c20-s8:eq:palabra-soporte-total-u031}
\end{equation}
et toutes ses coordonnées sont non nulles.

Pour chaque position \(b\), nous fixons
\begin{equation}
 P_b(c_*)=
 \begin{cases}
  \mathsf R,&(c_*)_b=1,\\
  I,&(c_*)_b=2,
 \end{cases}
 \qquad
 g_c:=\bigoplus_{b=0}^{11}P_b(c_*)CP_b(c_*)^{-1}.
 \label{p12-83-c20-s8:eq:isometria-orden-tres-u031}
\end{equation}

\begin{proposition}[Ordre, absence de points fixes et recollement]
\label{p12-83-c20-s8:prop:isometria-fijo-libre-pegado-u031}
L'opérateur \(g_c\) est d'ordre trois, n'a pas de vecteurs fixes non nuls et
préserve le réseau de Niemeier \(N(A_2^{12})\) construit dans U021.
\end{proposition}

\begin{proof}
Chaque bloc est conjugué au Coxeter \(C\), dont le polynôme minimal est
\(X^2+X+1\). Il est donc d'ordre trois et ne possède pas la valeur propre un. La somme
directe conserve les deux propriétés. Le Coxeter agit trivialement sur
\(A_2^*/A_2\) ; c'est pourquoi chaque bloc conserve la classe discriminante et la
réunion de classes latérales qui définit le recollement.
\end{proof}

\subsection{Préservation du voisin de Leech}

Écrivons
\[
 v:=v_{\mathcal F_*}
 =4\rho_1+\rho_2+\cdots+\rho_{12},
 \qquad
 v^2=54,
\]
pour le vecteur sélectionné par l'origine d'incidence de
\eqref{p12-83-c20-s5:eq:vector-marco-incidencial}. Soient
\[
 N:=N(A_2^{12}),
 \qquad
 N_0:=\ker(\langle\,\cdot\,,v\rangle\bmod3),
 \qquad
 N_v:=N_0+\mathbb Z\frac v3\cong\Lambda_{24}.
\]

\begin{theorem}[Préservation orientée du voisin]
\label{p12-83-c20-s8:thm:preservacion-vecino-orden-tres-u031}
Les déplacements
\begin{equation}
 y_*:=\frac{g_cv-v}{3},
 \qquad
 z_*:=\frac{g_c^{-1}v-v}{3}
 \label{p12-83-c20-s8:eq:desplazamientos-vecino-u031}
\end{equation}
appartiennent à \(N_0\). En conséquence,
\[
 g_cN_0=N_0,
 \qquad
 g_cN_v=N_v.
\]
\end{theorem}

\begin{proof}
Les mots discriminants de \(y_*\) et \(z_*\) sont \(c_*\) et
\(-c_*\), tous deux dans \(C_W\) ; d'où \(y_*,z_*\in N\). Le calcul dans chaque
facteur \(A_2\) donne
\begin{equation}
 \langle y_*,v\rangle
 =\langle z_*,v\rangle
 =-(4^2+11)=-27\equiv0\pmod3.
 \label{p12-83-c20-s8:eq:emparejamientos-desplazamientos-u031}
\end{equation}
Ils sont donc tous deux dans \(N_0\). Pour \(x\in N_0\),
\[
 \langle g_cx,v\rangle
 =\langle x,g_c^{-1}v\rangle
 =\langle x,v\rangle+3\langle x,z_*\rangle
 \equiv0\pmod3.
\]
Le même calcul avec \(g_c^{-1}\) donne l'égalité du noyau. Enfin,
\[
 g_c(v/3)=v/3+y_*,
\]
de sorte que la classe qui engendre le voisin est préservée.
\end{proof}

\subsection{\(24\) orientations et naturalité des réseaux}

L’énumérateur de poids de \eqref{p12-83-c20-s1:eq:enumerador-pesos-cw} contient
\(24y^{12}\). Par conséquent, l’ensemble
\[
 C_W^\times:=
 \{u\in C_W:u_b\in\{1,2\}\text{ pour tout }b\}
\]
possède 24 éléments. Pour chaque \(u\in C_W^\times\), nous définissons
\[
 P_b(u)=
 \begin{cases}
  \mathsf R,&u_b=1,\\
  I,&u_b=2,
 \end{cases}
 \qquad
 g_u:=\bigoplus_{b=0}^{11}P_b(u)CP_b(u)^{-1}.
\]

\begin{proposition}[Robustesse des orientations à support total]
\label{p12-83-c20-s8:prop:robustez-orientaciones-u031}
Pour toute \(u\in C_W^\times\),
\[
 \frac{g_uv-v}{3},
 \quad
 \frac{g_u^{-1}v-v}{3}
 \in N_0.
\]
En particulier, chaque \(g_u\) préserve le voisin construit avec
l’orientation transportée.
\end{proposition}

\begin{proof}
En coordonnées de racines simples,
\[
 C\rho-\rho=(-2,-1),
 \qquad
 C^{-1}\rho-\rho=(-1,-2).
\]
Les classes discriminantes des deux déplacements sont \(u\) et \(-u\).
Pour le coefficient radial \(a_b\in\{4,1\}\),
\[
 \left\langle
 \frac{(C^{\pm1}-I)a_b\rho}{3},a_b\rho
 \right\rangle=-a_b^2.
\]
La somme sur les douze facteurs donne de nouveau \(-27\). La preuve du
\cref{p12-83-c20-s8:thm:preservacion-vecino-orden-tres-u031} s’applique composante par
composante.
\end{proof}

Pour typer le recalibrage, soit
\(h\in\operatorname{Aut}(C_W)\) monomial : une permutation \(\sigma\) des
douze coordonnées suivie de multiplicateurs
\(\epsilon_b\in\{1,2\}\). Nous définissons son relèvement réticulaire
\begin{equation}
 \widetilde h
 :=P_\sigma\circ
 \bigoplus_{b=0}^{11}\mathsf R^{\nu_b},
 \qquad
 \nu_b=
 \begin{cases}
 0,&\epsilon_b=1,\\
 1,&\epsilon_b=2.
 \end{cases}
 \label{p12-83-c20-s8:eq:elevacion-monomial-u031}
\end{equation}
Son action sur le discriminant est exactement \(h\). Si le changement transporte
simultanément le code, le drapeau, l’origine incidencielle, le vecteur \(v\), le voisin et les
repères, alors
\begin{equation}
 g_{hu}=\widetilde h\,g_u\,\widetilde h^{-1}.
 \label{p12-83-c20-s8:eq:naturalidad-isometrias-u031}
\end{equation}
La matrice ternaire \(h\) n’agit sans élévation ni sur \(A_2^{12}\) ni sur
Leech.

\subsection{Trace de réseau et calcul du type \(0\)}

Soit \(\widehat g_c\) le relèvement standard de \(g_c\) à la VOA de réseau
\[
 V_{N_v}=M(1)\otimes\mathbb C_\varepsilon[N_v].
\]
Pour \(j=1,2\), la décomposition oscillateur--réseau produit
\begin{equation}
 \begin{split}
 Z_{0,j}(\tau)
 &:=
 \operatorname{Tr}_{V_{N_v}}
 \left(\widehat g_c^{\,j}q^{L_0-1}\right)\\
 &=q^{-1}
 \left(
 \sum_{\lambda\in N_v^{g_c^j}}
 \varepsilon_j(\lambda)
 q^{\langle\lambda,\lambda\rangle/2}
 \right)
 \prod_{n\ge1}\det(I-g_c^jq^n)^{-1}.
 \end{split}
 \label{p12-83-c20-s8:eq:factorizacion-traza-reticular-u031}
\end{equation}

\begin{proposition}[Trace des secteurs non triviaux]
\label{p12-83-c20-s8:prop:traza-reticular-t3-u031}
On a
\[
 Z_{0,1}(\tau)=Z_{0,2}(\tau)=t_3(\tau).
\]
\end{proposition}

\begin{proof}
L'absence de points fixes implique
\[
 N_v^{g_c}=N_v^{g_c^2}=\{0\},
\]
de sorte que la somme de réseau de
\eqref{p12-83-c20-s8:eq:factorizacion-traza-reticular-u031} vaut un. Dans chaque facteur,
\(g_c^j\) a pour valeurs propres \(\zeta_3,\zeta_3^2\) ; par conséquent,
\[
 \det(I-g_c^jq^n)=1+q^n+q^{2n}.
\]
Les douze facteurs reproduisent \eqref{p12-83-c20-s8:eq:t3-producto-u031}.
\end{proof}

\begin{theorem}[Poids tordu et type du relèvement]
\label{p12-83-c20-s8:thm:peso-torcido-tipo-cero-u031}
Le relèvement standard est d'ordre trois, de poids tordu
\begin{equation}
 \rho_{g_c}=\frac43
 \label{p12-83-c20-s8:eq:peso-torcido-u031}
\end{equation}
et de type zéro.
\end{theorem}

\begin{proof}
Dans
\(\mathfrak h=N_v\otimes_{\mathbb Z}\mathbb C\), chacun des douze
facteurs apporte un vecteur propre de valeur propre \(\zeta_3\) et un autre de valeur propre
\(\zeta_3^2\). Par conséquent,
\[
 \dim\mathfrak h_{\zeta_3}
 =\dim\mathfrak h_{\zeta_3^2}=12.
\]
La formule du poids du vide tordu d'ordre trois donne
\[
 \rho_{g_c}
 =\frac1{4\cdot3^2}
 \bigl(1\cdot2\cdot12+2\cdot1\cdot12\bigr)
 =\frac43.
\]
Pour une isométrie de réseau d'ordre impair, le relèvement standard conserve
l'ordre. Le type est la classe de
\[
 3^2\rho_{g_c}=12\equiv0\pmod3.
\]
Ainsi, les deux hypothèses initialement indépendantes sont calculées ici.
\end{proof}

\subsection{Orbifold cyclique et classe duale \texorpdfstring{\(3B\)}{3B}}

\begin{theorem}[Orbifold triadique du voisin de Leech]
\label{p12-83-c20-s8:thm:orbifold-triadico-vnatural-u031}
Les théorèmes classiques d'orbifold cyclique, appliqués au relèvement standard
du \cref{p12-83-c20-s8:thm:peso-torcido-tipo-cero-u031}, donnent
\begin{equation}
 \operatorname{Orb}_{\mathbb Z/3}
 (V_{N_v},\widehat g_c)
 \cong V^\natural.
 \label{p12-83-c20-s8:eq:orbifold-triadico-vnatural-u031}
\end{equation}
L'automorphisme dual appartient à la classe \(3B\) et sa série est
\begin{equation}
 T_{3B}(\tau)=t_3(\tau)+12.
 \label{p12-83-c20-s8:eq:serie-3b-u031}
\end{equation}
\end{theorem}

\begin{proof}
L'isométrie est sans points fixes et d'ordre trois ; le poids tordu et le type zéro ont été
calculés, non supposés. Les théorèmes d'orbifold cyclique identifient
l'extension holomorphe au module lunaire et classent l'automorphisme dual
comme \(3B\)
\cite{ChenLamShimakura2016,AbeLamYamada2017,Carnahan2017}.
L'identification utilise ces théorèmes après la construction de l'objet de réseau ;
elle ne se déduit pas uniquement d'une égalité de séries.
\end{proof}

Si \(V^{(i)}\) est le module \(\widehat g_c^{\,i}\)-tordu et
\[
 Z_{i,j}(\tau):=
 \operatorname{Tr}_{V^{(i)}}\!
 \left(
 \phi_i(\widehat g_c)^j q^{L_0-1}
 \right),
\]
le caractère de l'orbifold est la moyenne des neuf secteurs :
\begin{equation}
 \operatorname{ch}_{V^\natural}(\tau)
 =\frac13\sum_{i,j\in\mathbb Z/3}Z_{i,j}(\tau)
 =J(\tau).
 \label{p12-83-c20-s8:eq:promedio-nueve-sectores-u031}
\end{equation}
L'identité rationnelle de niveau trois est
\begin{equation}
 J
 =t_3+12+\frac{10\,3^9}{t_3}
 +\frac{4\,3^{14}}{t_3^2}
 +\frac{3^{18}}{t_3^3}.
 \label{p12-83-c20-s8:eq:j-desde-t3-u031}
\end{equation}
Comme \(t_3^{-1}=q+O(q^2)\),
\begin{equation}
 [q]J=54+10\cdot3^9=196884
 =54(5\cdot729+1).
 \label{p12-83-c20-s8:eq:coeficiente-196884-u031}
\end{equation}
C'est une compatibilité graduée au sein du corridor construit, non un
morphisme scalaire \(54\to V^\natural\).

\subsection{Construction de réseau d'ordre \(2\)}

La route d'ordre trois précédente et la construction classique de
Frenkel--Lepowsky--Meurman produisent deux présentations de la même VOA lunaire.
Pour déclarer la seconde, nous fixons l'extension centrale
\begin{equation}
 1\longrightarrow\{\pm1\}
 \longrightarrow\widehat\Lambda_{24}
 \longrightarrow\Lambda_{24}
 \longrightarrow1
 \label{p12-83-c20-s8:eq:extension-central-leech-u031}
\end{equation}
de commutateur
\((-1)^{\langle\lambda,\mu\rangle}\), et un cocycle
\(\varepsilon(\lambda,\mu)\) qui la représente. L'algèbre de groupe
tordue est \(\mathbb C_\varepsilon[\Lambda_{24}]\), et
\begin{equation}
 V_\Lambda
 :=M(1)\otimes\mathbb C_\varepsilon[\Lambda_{24}]
 \label{p12-83-c20-s8:eq:voa-reticular-leech-u031}
\end{equation}
est la VOA de réseau de charge centrale 24.

La négation \(\lambda\mapsto-\lambda\) possède un relèvement
\(\theta\) compatible avec \eqref{p12-83-c20-s8:eq:extension-central-leech-u031}. Sa trace
est
\begin{equation}
 t_2(\tau)
 :=
 \operatorname{Tr}_{V_\Lambda}
 (\theta q^{L_0-1})
 =q^{-1}\prod_{n\ge1}(1+q^n)^{-24}.
 \label{p12-83-c20-s8:eq:t2-leech-u031}
\end{equation}
En effet, la négation apparie \(e^\lambda\) avec \(e^{-\lambda}\), et
l'absence de torsion ne laisse que \(\lambda=0\) dans la trace de réseau ; chacun
des 24 oscillateurs apporte \((1+q^n)^{-1}\).

Soit \(V_\Lambda^T\) le module irréductible tordu par \(\theta\), et
\((V_\Lambda^T)^+\) son secteur de signe positif. La construction FLM est
\begin{equation}
 V^\natural
 =(V_\Lambda)^+\oplus(V_\Lambda^T)^+.
 \label{p12-83-c20-s8:eq:vnatural-flm-u031}
\end{equation}
La route \(\mathbb Z/2\) de
\eqref{p12-83-c20-s8:eq:vnatural-flm-u031} et la route \(\mathbb Z/3\) de
\eqref{p12-83-c20-s8:eq:orbifold-triadico-vnatural-u031} sont des constructions parallèles
identifiées par des théorèmes classiques ; aucune ne se substitue à l'autre.

\subsection{Groupe d'automorphismes, caractère et séries graduées}

Une fois \(V^\natural\) construit, deux publications distinctes apparaissent :
\begin{align}
 \operatorname{Aut}(V^\natural)&\cong\mathbb M,
 \label{p12-83-c20-s8:eq:aut-vnatural-monstruo-u031}\\
 \operatorname{ch}_{V^\natural}(\tau)
 &=J(\tau)=j(\tau)-744
 =q^{-1}+196884q+21493760q^2+\cdots.
 \label{p12-83-c20-s8:eq:caracter-vnatural-j-u031}
\end{align}
La première identifie le groupe de symétries ; la seconde compte les dimensions
graduées. La droite de vide est
\[
 (V^\natural)_0=\mathbb C\mathbf1,
\]
et l'absence de racines de \(\Lambda_{24}\) donne
\[
 (V^\natural)_1=0.
\]
En degré deux,
\[
 196884=1+196883=196560+324=54(5\cdot729+1).
\]
Seule la première égalité est une décomposition en dimensions de
représentations du Monstre ; les autres sont des contrôles arithmétiques.

Pour \(g\in\mathbb M\), on définit la série de McKay--Thompson
\begin{equation}
 T_g(\tau):=
 \operatorname{Tr}_{V^\natural}
 \left(gq^{L_0-1}\right).
 \label{p12-83-c20-s8:eq:serie-mckay-thompson-u031}
\end{equation}

\begin{theorem}[Moonshine monstrueux]
\label{p12-83-c20-s8:thm:moonshine-monstruoso-u031}
L'action graduée de \(\mathbb M\) sur \(V^\natural\) produit la famille
\(\{T_g\}_{g\in\mathbb M}\). Les théorèmes de Moonshine identifient ces
séries à des fonctions modulaires de genre zéro et prouvent leurs identités de
réplicabilité
\cite{ConwayNorton1979,Borcherds1992,FLM1988}.
\end{theorem}

Ce théorème a pour entrée le couple
\[
 \bigl(V^\natural,\operatorname{Aut}(V^\natural)\bigr)
\]
et sa graduation. Ce n'est pas la conséquence d'une flèche dont le domaine serait le
groupe abstrait \(\mathbb M\).

Pour \(J=T_{1A}\), soient \(P_n\) les polynômes de Faber normalisés par
\[
 P_n(J)=q^{-n}+O(q).
\]
L'identité de Hecke--Faber est
\begin{equation}
 P_n(J)=n\,T_nJ.
 \label{p12-83-c20-s8:eq:hecke-faber-u031}
\end{equation}
Elle contrôle tous les degrés ; le coefficient 196884 n'est que la première
manifestation non triviale.

\subsection{Graphe de dépendances du corridor Paley–Witt–Golay–Leech–Moonshine}

\begin{figure}[htbp]
\centering
\begin{tikzcd}[column sep=huge,row sep=large]
 &C_W
   \arrow[dl,"W_{12}"']
   \arrow[d,"N(A_2^{12})"',font=\scriptsize]
   \arrow[dr,phantom,"\substack{\text{branches}\\\text{distinctes}}",
          description,pos=.72,font=\scriptsize]&\\
 M_{12}
 &\Lambda_{24}
   \arrow[d,"V_\Lambda"]
 &(\mathcal G_{24},D,\ell)
   \arrow[d,"W_{24}"]\\
 &V^\natural
   \arrow[dl,"\operatorname{Aut}"']
   \arrow[d,"\operatorname{ch}"]
   \arrow[dr,"g\mapsto T_g"]
 &M_{24}\\
 \mathbb M
 &J
 &\{T_g\}_{g\in\mathbb M}
\end{tikzcd}
\caption{Bifurcation exceptionnelle et trois publications de
\(V^\natural\). L'entrée binaire \((\mathcal G_{24},D,\ell)\) est
indépendante ; il n'existe pas de flèche \(W_{24}\to\Lambda_{24}\).}
\label{p12-83-c20-s8:fig:grafo-excepcional-moonshine-u031}
\end{figure}

\begin{criterion}[Critères de réfutation du corridor lunaire]
\label{p12-83-c20-s8:crit:refutacion-corredor-lunar-u031}
La construction est réfutée si l'un des faits suivants se produit :
\begin{enumerate}
 \item \(\lambda\) n'est pas bijective ou
       \eqref{p12-83-c20-s8:eq:identidad-entrelazamiento-app-witt-u031} échoue ;
 \item \(c_*\notin C_W\), \(g_c\) possède des points fixes ou ne préserve pas le
       voisin ;
 \item les déplacements de
       \eqref{p12-83-c20-s8:eq:desplazamientos-vecino-u031} n'appartiennent pas à \(N_0\) ;
 \item le relèvement de réseau \(\widetilde h\) ne réalise pas l'automorphisme
       monomial \(h\) dans le discriminant ;
 \item le poids tordu n'est pas \(4/3\) ou le relèvement n'est pas de type zéro ;
 \item l'orbifold est identifié à \(V^\natural\) sans vérification des
       hypothèses du théorème classique ;
 \item on utilise \(W_{24}\) comme constructeur de Leech ;
 \item une égalité de dimensions est présentée comme construction de
       \(\mathbb M\) ;
 \item on sérialise indûment les publications
       \(\operatorname{Aut}\), \(\operatorname{ch}\) et \(g\mapsto T_g\).
\end{enumerate}
\end{criterion}
